\section{Локальная лемма Ловаса. Числа Рамсея}

\subsection{Локальная лемма Ловаса}

\subsubsection*{Напоминание о независимости}

Пусть даны события $A, B_1, \dots, B_n$.
Будем говорить, что $A$ независимо от семейства $B_1,\dots,B_n$, если для любого
$I \subseteq \{1,\dots,n\}$ выполнено
\[
\Pr\!\left(A \mid \bigwedge_{i \in I} B_i\right)=\Pr(A),
\]
если условная вероятность определена.

\begin{Definition}{Граф зависимостей}{}
	Пусть $\mathcal A=\{A_1,\dots,A_n\}$ — семейство событий.
	Ориентированный граф $\Gamma$ называется \emph{графом зависимостей} для $\mathcal A$,
	если $V(\Gamma)=\{1,\dots,n\}$ и для каждого $i$ событие $A_i$ независимо от всех событий
	\[
	\{A_k : k \notin \Gamma^+(i)\},
	\]
	где $\Gamma^+(i)$ — множество исходящих соседей вершины $i$.
\end{Definition}

\begin{Example}{Монеты}{}
	Пусть есть $n-1$ независимых честных монет. Обозначим:
	\[
	A_i=\{\text{$i$-я монета выпала орлом}\}, \qquad i=1,\dots,n-1,
	\]
	\[
	A_n=\{\text{число орлов чётно}\}.
	\]
	Тогда каждое событие независимо от любого набора из не более чем $n-2$ остальных событий,
	но не обязано быть независимо от всех остальных сразу.
	
	Поэтому для этого семейства графом зависимостей будет любой граф, в котором у каждой вершины
	есть хотя бы один исходящий сосед.
\end{Example}

\begin{Theorem}{Локальная лемма Ловаса, асимметричная форма}{} 
	Пусть $\mathcal A=\{A_1,\dots,A_n\}$ — семейство событий, а $\Gamma$ — граф зависимостей
	для $\mathcal A$.
	
	Пусть существует функция $f:\mathcal A \to [0,1)$ такая, что для каждого $i$ выполнено
	\[
	\Pr(A_i)
	\le
	f(A_i)\prod_{k\in \Gamma^+(i)}\bigl(1-f(A_k)\bigr).
	\]
	Тогда
	\[
	\Pr\!\left(\overline{A_1}\wedge \dots \wedge \overline{A_n}\right)
	\ge
	\prod_{i=1}^{n}\bigl(1-f(A_i)\bigr).
	\]
	В частности, эта вероятность положительна.
\end{Theorem}

\begin{proof}
	Докажем вспомогательное утверждение.
	
	\medskip
	\noindent
	\textbf{Утверждение.}
	Для любого $S\subseteq \mathcal A$ и любого $A_i\notin S$ выполнено:
	\[
	\Pr\!\left(\bigwedge_{B\in S}\overline B\right)>0 \text{ и }
	\Pr\!\left(A_i\mid \bigwedge_{B\in S}\overline B\right)\le f(A_i).
	\]
	
	\medskip
	Докажем утверждение индукцией по $|S|$.
	
	\medskip
	\noindent
	\textbf{База:} $S=\varnothing$.
	
	Тогда
	\[
	\Pr(A_i)
	\le
	f(A_i)\prod_{k\in \Gamma^+(i)}\bigl(1-f(A_k)\bigr)
	\le f(A_i),
	\]
	поскольку все множители в произведении не превосходят $1$.
	
	\medskip
	\noindent
	\textbf{Переход.}
	Пусть утверждение доказано для всех множеств размера меньше $|S|$.
	Разобьём $S$ на две части:
	\[
	S_1=\{B\in S : B=A_k,\ k\in \Gamma^+(i)\},
	\]
	\[
	S_2=S\setminus S_1.
	\]
	То есть $S_1$ состоит из событий, которые могут зависеть с $A_i$, а $S_2$ — из событий,
	от которых $A_i$ независимо.
	
	Сначала докажем положительность вероятности.
	Возьмём любое $A_k\in S$. Тогда
	\[
	\Pr\!\left(\bigwedge_{B\in S}\overline B\right)
	=
	\Pr\!\left(
	\overline{A_k}
	\mid
	\bigwedge_{B\in S\setminus\{A_k\}}\overline B
	\right)
	\Pr\!\left(
	\bigwedge_{B\in S\setminus\{A_k\}}\overline B
	\right).
	\]
	По предположению индукции
	\[
	\Pr\!\left(
	A_k
	\mid
	\bigwedge_{B\in S\setminus\{A_k\}}\overline B
	\right)
	\le f(A_k),
	\]
	следовательно,
	\[
	\Pr\!\left(
	\overline{A_k}
	\mid
	\bigwedge_{B\in S\setminus\{A_k\}}\overline B
	\right)
	\ge 1-f(A_k)>0.
	\]
	Повторяя это для элементов множества $S$, получаем
	\[
	\Pr\!\left(\bigwedge_{B\in S}\overline B\right)
	\ge
	\prod_{B\in S}\bigl(1-f(B)\bigr)>0.
	\]
	
	Теперь оценим условную вероятность.
	Пусть $S_1=\{B_1,\dots,B_\ell\}.$ Тогда
	\[
	\Pr\!\left(A_i\mid \bigwedge_{B\in S}\overline B\right)
	=
	\frac{
		\Pr\!\left(
		A_i\wedge \bigwedge_{B\in S_1}\overline B
		\mid
		\bigwedge_{B\in S_2}\overline B
		\right)
	}{
		\Pr\!\left(
		\bigwedge_{B\in S_1}\overline B
		\mid
		\bigwedge_{B\in S_2}\overline B
		\right)
	}.
	\]
	
	Оценим числитель:
	\[
	\Pr\!\left(
	A_i\wedge \bigwedge_{B\in S_1}\overline B
	\mid
	\bigwedge_{B\in S_2}\overline B
	\right)
	\le
	\Pr\!\left(
	A_i
	\mid
	\bigwedge_{B\in S_2}\overline B
	\right).
	\]
	Так как все события из $S_2$ не являются исходящими соседями $A_i$, то по определению
	графа зависимостей
	\[
	\Pr\!\left(
	A_i
	\mid
	\bigwedge_{B\in S_2}\overline B
	\right)
	=
	\Pr(A_i).
	\]
	Значит,
	\[
	\Pr\!\left(
	A_i\wedge \bigwedge_{B\in S_1}\overline B
	\mid
	\bigwedge_{B\in S_2}\overline B
	\right)
	\le
	f(A_i)\prod_{k\in \Gamma^+(i)}\bigl(1-f(A_k)\bigr).
	\]
	
	Оценим знаменатель. По цепному правилу:
	\[
	\Pr\!\left(
	\bigwedge_{r=1}^{\ell}\overline{B_r}
	\mid
	\bigwedge_{B\in S_2}\overline B
	\right)
	=
	\prod_{r=1}^{\ell}
	\Pr\!\left(
	\overline{B_r}
	\mid
	\bigwedge_{B\in S_2}\overline B
	\wedge
	\bigwedge_{t<r}\overline{B_t}
	\right).
	\]
	В каждом множителе условие содержит меньше чем $|S|$ событий, поэтому по предположению
	индукции
	\[
	\Pr\!\left(
	\overline{B_r}
	\mid
	\bigwedge_{B\in S_2}\overline B
	\wedge
	\bigwedge_{t<r}\overline{B_t}
	\right)
	\ge
	1-f(B_r).
	\]
	Следовательно,
	\[
	\Pr\!\left(
	\bigwedge_{B\in S_1}\overline B
	\mid
	\bigwedge_{B\in S_2}\overline B
	\right)
	\ge
	\prod_{B\in S_1}\bigl(1-f(B)\bigr).
	\]
	
	Итак,
	\[
	\Pr\!\left(A_i\mid \bigwedge_{B\in S}\overline B\right)
	\le
	\frac{
		f(A_i)\prod_{k\in \Gamma^+(i)}\bigl(1-f(A_k)\bigr)
	}{
		\prod_{B\in S_1}\bigl(1-f(B)\bigr)
	}.
	\]
	Так как $S_1\subseteq \{A_k:k\in\Gamma^+(i)\}$, то часть множителей сокращается, а остальные
	не превосходят $1$. Значит,
	\[
	\Pr\!\left(A_i\mid \bigwedge_{B\in S}\overline B\right)
	\le f(A_i).
	\]
	Вспомогательное утверждение доказано.
	
	\medskip
	Теперь применим его к исходной задаче:
	\[
	\Pr\!\left(\overline{A_1}\wedge\dots\wedge\overline{A_n}\right)
	=
	\prod_{i=1}^{n}
	\Pr\!\left(
	\overline{A_i}
	\mid
	\overline{A_{i+1}}\wedge\dots\wedge\overline{A_n}
	\right).
	\]
	Каждый множитель равен $1-
	\Pr\!\left(
	A_i
	\mid
	\overline{A_{i+1}}\wedge\dots\wedge\overline{A_n}
	\right)
	\ge 1-f(A_i).$
	
	Следовательно, $\Pr\!\left(\overline{A_1}\wedge\dots\wedge\overline{A_n}\right)
	\ge
	\prod_{i=1}^{n}\bigl(1-f(A_i)\bigr).$
\end{proof}

\begin{Theorem}{Локальная лемма Ловаса, симметричная форма}{}
	Пусть $\mathcal A=\{A_1,\dots,A_n\}$ — семейство событий.
	Пусть
	\[
	\Pr(A_i)\le q
	\]
	для всех $i$, и пусть существует граф зависимостей $\Gamma$, в котором исходящая степень
	каждой вершины не превосходит $d$.
	
	Если
	\[
	\mathrm e\,q(d+1)\le 1,
	\]
	то
	\[
	\Pr\!\left(\overline{A_1}\wedge\dots\wedge\overline{A_n}\right)>0.
	\]
\end{Theorem}

\begin{proof}
	Положим
	\[
	f(A_i)=\frac{1}{d+1}
	\]
	для всех $i$.
	
	Тогда
	\[
	f(A_i)\prod_{k\in\Gamma^+(i)}\bigl(1-f(A_k)\bigr)
	\ge
	\frac{1}{d+1}
	\left(1-\frac{1}{d+1}\right)^d.
	\]
	Имеем
	\[
	\frac{1}{d+1}
	\left(1-\frac{1}{d+1}\right)^d
	=
	\frac{1}{d+1}
	\left(\frac d{d+1}\right)^d
	=
	\frac{1}{d+1}
	\cdot
	\frac{1}{\left(1+\frac1d\right)^d}
	\ge
	\frac{1}{\mathrm e(d+1)}.
	\]
	Из условия $\mathrm e q(d+1)\le 1$ следует
	\[
	q\le \frac{1}{\mathrm e(d+1)}.
	\]
	Значит,
	\[
	\Pr(A_i)\le q
	\le
	f(A_i)\prod_{k\in\Gamma^+(i)}\bigl(1-f(A_k)\bigr).
	\]
	Условия асимметричной формы ЛЛЛ выполнены.
\end{proof}

\begin{Corollary}{Удобная грубая форма}{}
	Пусть $\Pr(A_i)\le q$, и каждое событие зависит не более чем от $d$ других событий.
	Если
	\[
	4qd\le 1,
	\]
	то
	\[
	\Pr\!\left(\overline{A_1}\wedge\dots\wedge\overline{A_n}\right)>0.
	\]
\end{Corollary}


\begin{Example}{$k$-КНФ}{}
Рассмотрим $k$-КНФ:
\[
C_1\wedge C_2\wedge\dots\wedge C_m,
\]
где каждый дизъюнкт $C_i$ содержит ровно $k$ литералов.

Выберем случайное значение каждой переменной независимо и равновероятно.
Обозначим
\[
A_i=\{\text{дизъюнкт $C_i$ не выполнен}\}.
\]
Тогда
\[
\Pr(A_i)=2^{-k}.
\]

События $A_i$ и $A_j$ зависимы только тогда, когда дизъюнкты $C_i$ и $C_j$ имеют общую
переменную.

Если каждый дизъюнкт имеет общие переменные не более чем с $2^{k-2}$ другими дизъюнктами,
то можно взять
\[
q=2^{-k},
\qquad
d\le 2^{k-2}.
\]
Тогда
\[
4qd
\le
4\cdot 2^{-k}\cdot 2^{k-2}
=1.
\]
По локальной лемме Ловаса
\[
\Pr\!\left(\overline{A_1}\wedge\dots\wedge\overline{A_m}\right)>0.
\]
Значит, существует набор значений переменных, при котором все дизъюнкты выполнены.
То есть такая $k$-КНФ выполнима.
\end{Example}

\begin{Example}{независимое множество из списков}{}
Пусть $G$ — граф с максимальной степенью не больше $d$.
Пусть даны попарно непересекающиеся множества вершин
\[
V_1,V_2,\dots,V_t,
\]
причём $|V_i|\ge 2\mathrm e d$
для всех $i$.

Докажем, что можно выбрать по одной вершине $u_i\in V_i$
так, чтобы множество $\{u_1,\dots,u_t\}$
было независимым.

Положим
\[
K=\lfloor 2\mathrm e d\rfloor+1.
\]
Можно уменьшить каждое множество $V_i$ до размера $K$.

Теперь из каждого $V_i$ независимо и равновероятно выберем одну вершину.
Для каждого ребра $r$ графа $G$, концы которого лежат в разных множествах $V_i$, определим
событие
\[
A_r=\{\text{выбраны оба конца ребра $r$}\}.
\]
Тогда
\[
\Pr(A_r)=\frac{1}{K^2}.
\]

Событие $A_r$ зависит только от событий, связанных с теми же двумя множествами $V_i$,
из которых берутся концы ребра $r$.

Если $r$ соединяет вершины из $V_i$ и $V_j$, то число рёбер, которые имеют хотя бы один конец
в $V_i\cup V_j$, не превосходит
\[
2Kd.
\]
Исключая само ребро $r$, получаем оценку на число зависимостей:
\[
D\le 2(Kd-1).
\]

Проверим условие симметричной ЛЛЛ:
\[
\mathrm e\cdot \frac{1}{K^2}\cdot (D+1)
\le
\mathrm e\cdot \frac{1}{K^2}\cdot \bigl(2(Kd-1)+1\bigr).
\]
Так как $K>2\mathrm e d$, то
\[
\mathrm e\cdot \frac{2(Kd-1)+1}{K^2}
<
\frac{2\mathrm e d}{K}
<1.
\]
Следовательно, с положительной вероятностью не произойдёт ни одно событие $A_r$.

Это означает, что выбранные вершины не соединены рёбрами.
Значит, существует независимое множество, содержащее ровно одну вершину из каждого $V_i$.
\end{Example}

\subsection{Числа Рамсея}

\begin{Definition}{Числа Рамсея}{}
	$R(m,n)$ — наименьшее натуральное число $N$ такое, что в любом простом графе на $N$
	вершинах содержится либо полный подграф на $m$ вершинах, либо независимое множество
	на $n$ вершинах.
	
	Эквивалентная формулировка: если рёбра полного графа $K_N$ покрасить в два цвета,
	то найдётся либо полный подграф $K_m$ первого цвета, либо полный подграф $K_n$ второго цвета.
\end{Definition}

\begin{Proposition}{}{}
	Для чисел Рамсея выполнено:
	\[
	R(1,n)=R(m,1)=1,
	\]
	\[
	R(m,n)=R(n,m).
	\]
\end{Proposition}

\begin{Theorem}{Рекуррентная верхняя оценка}{}
	Для всех $m,n\ge 2$ выполнено
	\[
	R(m,n)\le R(m-1,n)+R(m,n-1).
	\]
\end{Theorem}

\begin{proof}
	Обозначим
	\[
	N=R(m-1,n)+R(m,n-1).
	\]
	Рассмотрим любую раскраску рёбер полного графа $K_N$ в красный и белый цвета.
	
	Выберем произвольную вершину $v$.
	Остальные $N-1$ вершин разобьём на два множества:
	\[
	X=\{u : vu \text{ — красное ребро}\},
	\]
	\[
	Y=\{u : vu \text{ — белое ребро}\}.
	\]
	
	Если
	\[
	|X|\ge R(m-1,n),
	\]
	то внутри $X$ найдётся либо красный $K_{m-1}$, либо белый $K_n$.
	В первом случае вместе с вершиной $v$ получаем красный $K_m$.
	Во втором случае уже есть белый $K_n$.
	
	Если
	\[
	|Y|\ge R(m,n-1),
	\]
	то внутри $Y$ найдётся либо красный $K_m$, либо белый $K_{n-1}$.
	В первом случае уже есть красный $K_m$.
	Во втором случае вместе с вершиной $v$ получаем белый $K_n$.
	
	Остаётся заметить, что не могут одновременно выполняться неравенства
	\[
	|X|\le R(m-1,n)-1,
	\qquad
	|Y|\le R(m,n-1)-1,
	\]
	потому что тогда
	\[
	N-1=|X|+|Y|
	\le
	R(m-1,n)+R(m,n-1)-2
	=
	N-2,
	\]
	что невозможно.
\end{proof}

\begin{Corollary}{}{}
	Для всех $m,n\ge 1$ выполнено
	\[
	R(m,n)\le \binom{m+n-2}{m-1}
	=
	\binom{m+n-2}{n-1}.
	\]
\end{Corollary}

\begin{Remark}{}{}
	Локальная лемма Ловаса часто даёт хорошие нижние оценки на числа Рамсея, особенно для
	$R(\ell,k)$ при малых фиксированных $\ell$.
\end{Remark}

\begin{Example}{нижняя оценка на $R(3,k)$}{}
Рассмотрим случайный граф $G$ на $N$ вершинах.
Каждое ребро проводим независимо с вероятностью $q$.

Для каждого $S\subseteq V(G)$, $|S|=3$, определим событие
\[
A_S=\{\text{на $S$ есть полный граф } K_3\}.
\]
Для каждого $T\subseteq V(G)$, $|T|=k$, определим событие
\[
B_T=\{\text{на $T$ пустой граф}\}.
\]

Тогда
\[
\Pr(A_S)=q^3,
\]
поскольку треугольник содержит $3$ ребра, и
\[
\Pr(B_T)=(1-q)^{\binom{k}{2}},
\]
поскольку в $T$ есть $\binom{k}{2}$ возможных рёбер, и все они должны отсутствовать.

События зависят только тогда, когда соответствующие множества вершин имеют общую пару
вершин. Поэтому можно использовать следующие оценки на число зависимостей:
\[
A_S \text{ зависит не более чем от } 3(N-3) \text{ событий вида } A_{S'},
\]
\[
A_S \text{ зависит не более чем от } 3\binom{N-2}{k-2} \text{ событий вида } B_T,
\]
\[
B_T \text{ зависит не более чем от } \binom{k}{2}(N-2) \text{ событий вида } A_S,
\]
\[
B_T \text{ зависит не более чем от } \binom{k}{2}\binom{N-2}{k-2}
\text{ событий вида } B_{T'}.
\]

Применим асимметричную форму ЛЛЛ.
Положим
\[
f(A_S)=x,
\qquad
f(B_T)=y.
\]
Тогда достаточно подобрать $x,y\in[0,1)$ так, чтобы выполнялись неравенства
\[
q^3
\le
x(1-x)^{3(N-3)}
(1-y)^{3\binom{N-2}{k-2}},
\]
\[
(1-q)^{\binom{k}{2}}
\le
y(1-x)^{\binom{k}{2}(N-2)}
(1-y)^{\binom{k}{2}\binom{N-2}{k-2}}.
\]

Если такие параметры существуют, то с положительной вероятностью в графе нет ни треугольника,
ни независимого множества размера $k$.
Следовательно,
\[
R(3,k)>N.
\]

При подходящем выборе параметров этим методом получается оценка
\[
R(3,k)
\ge
\Omega\!\left(\left(\frac{k}{\log k}\right)^2\right).
\]
\end{Example}